% [joc_prep test copy generated by make_test_copies.py from ../main.tex -- do not edit]
% =====================================================================
%  supplementary_header_blind.tex -- JoC header for the supplementary file.
%  Replaces everything in supplementary.tex from \documentclass down to and
%  including \maketitle.  Referee layout (joc_referee.tex), JoC reference
%  style (joc_bib.tex) and the double-anonymous switch (blind.tex): in blind
%  mode the author/affiliation lines below are swallowed and replaced by one
%  neutral line; with \blindfalse they print as in the non-blind version.
%  The title deliberately does not name the authors or their institutions.
% =====================================================================
\documentclass[pdflatex,referee]{sn-jnl}
\input{joc_referee}
\input{joc_bib}
\input{blind}
\pdfmapfile{=local.map}          %% build-machine workaround; delete elsewhere
\usepackage{graphicx}%
\usepackage{multirow}%
\usepackage{amsmath,amssymb,amsfonts}%
\usepackage{amsthm}%
\usepackage{mathrsfs}%
\usepackage[title]{appendix}%
\usepackage{xcolor}%
\usepackage{textcomp}%
\usepackage{manyfoot}%
\usepackage{booktabs}%

\renewcommand{\thefigure}{S\arabic{figure}}
\renewcommand{\thetable}{S\arabic{table}}
\renewcommand{\thesection}{S\arabic{section}}

\begin{document}

\title{Supplementary Material\\
       \large Downstream Utility of Synthetic Data for Clustering}
\author[1]{Oriol Farr\'es Vilar}
\author[2]{Jordi Cort\'es}
\author[3]{Xinnuo Chen}
\author[4]{Victoria L\'opez}
\author[2]{Daniel Fern\'andez}
\affil[1]{Department of Computer Science,
          Universitat Polit\`ecnica de Catalunya, Barcelona, Spain}
\affil[2]{Department of Statistics and Operations Research,
          Universitat Polit\`ecnica de Catalunya, Barcelona, Spain}
\affil[3]{Department of Statistical Science,
          University College London, London, United Kingdom}
\affil[4]{CUNEF Universidad, Madrid, Spain}
\abstract{Supporting figures and tables for the main article
  \blindalt{(submitted for anonymous review to the Journal of Classification)}{}.
  Code: \repourl.}
\maketitle


\section{Recovery of the number of groups by cluster count and separation}

Figure~\ref{figS:heatmap} presents the recovery rate of the true number of
groups as a function of the true cluster count $k$ and the separation
$\sigma$, for both algorithms. It complements panel (a) of Figure~3 of
the main article in a different layout, making the $k \times \sigma$
interaction directly legible: recovery collapses in the lower-left quadrant
(many groups, little separation) and is essentially perfect in the upper-right.

\begin{figure}[h]
\centering
\includegraphics[width=\textwidth]{figures/figS_heatmap.pdf}
\caption{Recovery rate of the true number of groups by true cluster count and
  separation, for $k$-means (left) and Ward (right). Cell values are
  proportions over the full design.}
\label{figS:heatmap}
\end{figure}

\section{Agreement in the estimated number of groups between original and
         synthetic data}

Table~\ref{tabS:khat} cross-tabulates the number of groups estimated from the
original data against the number estimated from the corresponding synthetic
data, stratified by separation. Unlike the recovery rate in the main article,
which scores $\hat{k}$ against the construction-time ground truth, this table
compares the two estimates with each other and so isolates the effect of
synthesis from the difficulty of the clustering problem itself.

The pattern confirms the expectation that at high separation the mass
concentrates on the diagonal: exact agreement rises from $68.4\%$ at
$\sigma = 0.1$ to $98.2\%$ at $\sigma = 10$ for $k$-means (Ward: $64.6\%$ to
$97.5\%$). In the overlapping regime the off-diagonal mass is not symmetric.
At $\sigma = 0.1$ the synthetic data yield a \emph{lower} estimate than the
original in $15.2\%$ of comparisons and a higher one in $16.4\%$; the largest
single off-diagonal cell is the original data supporting $\hat{k} \ge 5$ while
the synthetic data support $\hat{k} = 2$ ($8.4\%$). That is, when the original
data contain weak, diffuse structure that the silhouette criterion resolves
into many small groups, CART synthesis tends to smooth it away into a single
coarse split. The tendency disappears once the groups are genuinely separated.

\begin{table}[h]
\centering
\caption{Agreement in the estimated number of groups between original and
  synthetic data, by separation. Rows are $\hat{k}$ from the original data,
  columns $\hat{k}$ from the synthetic data; cells are percentages of all
  comparisons at that separation (each block sums to $100\%$). Estimates of
  five or more groups are pooled for display, but the agreement column is
  computed on unpooled values. Both data sources are searched over a fixed
  range $\hat{k} \in \{2,\dots,8\}$, so the comparison is symmetric; under the
  narrower per-scenario range used for the recovery rates in the main text,
  agreement is higher in the overlapping regime ($75.1\%$ and $67.6\%$ at
  $\sigma = 0.1$, $79.7\%$ and $75.9\%$ at $\sigma = 2$, for $k$-means and Ward
  respectively) and essentially unchanged at $\sigma \ge 6$, where the two
  conventions differ by at most $0.1$ percentage points.
  $n = 180{,}000$ comparisons per separation level.}
\label{tabS:khat}
\small
\input{supp_table.tex}
\end{table}


\section{Sensitivity to the candidate range for the number of groups}

In the main article the silhouette criterion searches
$\hat{k} \in \{2,\dots,\min(k+3,8)\}$, a rule inherited from earlier work that
ties the search range to the planted number of groups. Table~\ref{tabS:range}
repeats the count analyses with the fixed range $\{2,\dots,8\}$, applied
identically to original and synthetic data. Widening the range lowers both
agreement measures in the overlapping regime, where the silhouette criterion
occasionally selects large values of $\hat{k}$, and leaves them essentially
unchanged at $\sigma \ge 6$. The qualitative conclusions are unaffected: the
two algorithms remain indistinguishable against the planted number of groups
and $k$-means remains significantly more stable against the original data.

\begin{table}[h]
\centering
\caption{Recovery of the planted number of groups (RR) and agreement with the
  number estimated from the original data (AG), by separation, under the
  per-scenario candidate range used in the main text and under a fixed range.
  Wilcoxon $p$: paired signed-rank test of $k$-means versus Ward across the 144
  scenario-level means.}
\label{tabS:range}
\small
\begin{tabular}{lrrrrrrrr}
\toprule
 & \multicolumn{4}{c}{$\hat{k} \in \{2,\dots,\min(k+3,8)\}$} & \multicolumn{4}{c}{$\hat{k} \in \{2,\dots,8\}$} \\
\cmidrule(lr){2-5}\cmidrule(lr){6-9}
 & \multicolumn{2}{c}{RR} & \multicolumn{2}{c}{AG} & \multicolumn{2}{c}{RR} & \multicolumn{2}{c}{AG} \\
\cmidrule(lr){2-3}\cmidrule(lr){4-5}\cmidrule(lr){6-7}\cmidrule(lr){8-9}
Separation & $k$-m. & Ward & $k$-m. & Ward & $k$-m. & Ward & $k$-m. & Ward \\
\midrule
$\sigma=0.1$ & 0.272 & 0.296 & 0.751 & 0.676 & 0.247 & 0.281 & 0.684 & 0.646 \\
$\sigma=2$ & 0.304 & 0.307 & 0.797 & 0.759 & 0.288 & 0.295 & 0.764 & 0.738 \\
$\sigma=6$ & 0.810 & 0.789 & 0.926 & 0.882 & 0.810 & 0.789 & 0.926 & 0.882 \\
$\sigma=10$ & 0.954 & 0.950 & 0.983 & 0.976 & 0.954 & 0.950 & 0.982 & 0.975 \\
\midrule
Overall & 0.585 & 0.585 & 0.864 & 0.823 & 0.575 & 0.579 & 0.839 & 0.811 \\
Wilcoxon $p$ & \multicolumn{2}{c}{0.39} & \multicolumn{2}{c}{$<0.001$} & \multicolumn{2}{c}{0.76} & \multicolumn{2}{c}{$<0.001$} \\
\bottomrule
\end{tabular}
\end{table}

\section{Independent-resampling baseline}

To judge whether a given level of agreement is high, it helps to know how
often two datasets that differ only by sampling variability yield the same
estimate. Within each scenario we compare the estimates $\hat{k}_R$ from the
$n = 5$ independent original datasets and record the proportion of the 10
pairs that coincide. Table~\ref{tabS:baseline} sets this baseline beside the
synthetic-versus-original agreement AG. Outside the most overlapping regime,
a CART synthetic dataset agrees with the original dataset from which it was
generated more often than two independent samples from the same population
agree with each other. At $\sigma = 0.1$ the two are comparable for Ward, and
synthetic agreement is lower for $k$-means. The baseline rests on only 10
pairs per scenario and is therefore a coarse reference rather than a formal
null distribution.

\begin{table}[h]
\centering
\caption{Synthetic-versus-original agreement (AG) and the
  independent-resampling baseline (proportion of pairs of independent original
  datasets in the same scenario with equal $\hat{k}_R$). Means over scenarios;
  candidate range $\{2,\dots,\min(k+3,8)\}$.}
\label{tabS:baseline}
\small
\begin{tabular}{llrrrr}
\toprule
 & & \multicolumn{2}{c}{$k$-means} & \multicolumn{2}{c}{Ward} \\
\cmidrule(lr){3-4}\cmidrule(lr){5-6}
Stratum & Level & AG & Baseline & AG & Baseline \\
\midrule
Separation $\sigma$ & $0.1$ & 0.751 & 0.817 & 0.676 & 0.681 \\
 & $2$ & 0.797 & 0.750 & 0.759 & 0.728 \\
 & $6$ & 0.926 & 0.764 & 0.882 & 0.758 \\
 & $10$ & 0.983 & 0.933 & 0.976 & 0.922 \\
\addlinespace
Distribution & Normal & 0.896 & 0.872 & 0.898 & 0.869 \\
 & Gamma & 0.833 & 0.760 & 0.749 & 0.675 \\
\midrule
Overall & & 0.864 & 0.816 & 0.823 & 0.772 \\
\bottomrule
\end{tabular}
\end{table}

\end{document}
